Los resultados caracterizan un protocolo con presupuesto compartido, inicialización aleatoria y una partición congelada. El número de parámetros y el diseño de la representación producen costos distintos, pero su relación con el F1 depende del entrenamiento observado. La pérdida, el F1 de validación y las predicciones de prueba permiten estudiar estas diferencias sin usar prueba para seleccionar hiperparámetros.

El resultado histórico del notebook de referencia se obtuvo con otra división y seleccionando por exactitud de validación. La retirada de duplicados y el agrupamiento de imágenes similares modifican la dificultad y el tamaño de las muestras. Por ello, el contraste con esa cifra exige considerar esas diferencias. El valor de la reproducción está en conservar la arquitectura y hacer trazables los cambios de protocolo.

La ausencia de identificadores de paciente es una limitación central. Un agrupamiento visual puede reducir fugas detectables y aun así dejar cortes relacionados en conjuntos distintos. La heterogeneidad de fuentes agrega otra incertidumbre. El formato y los rasgos de adquisición pueden estar asociados con las etiquetas, y los mapas de activación ofrecen una comprobación exploratoria que requiere una validación clínica y por fuente más detallada.

La validación reúne 526 imágenes y se consulta al final de cada época. Su tamaño y el agrupamiento visual pueden contribuir a oscilaciones del F1 y a cambios del checkpoint elegido. Las curvas se conservan para examinar esta sensibilidad. La comparación final mantiene el conjunto de prueba reservado durante todas esas decisiones.

El uso de cinco semillas permite cuantificar variación debida al entrenamiento, pero conserva el mismo conjunto de prueba. La prueba exacta bilateral tiene resolución limitada y requiere su hipótesis de intercambiabilidad de signos. La magnitud del efecto y las diferencias individuales aportan información adicional. La variabilidad entre centros y pacientes requeriría particiones por paciente verificables y validación externa.

El entrenamiento desde cero alinea la condición de partida. La receta de Adam compartida facilita comparar las arquitecturas bajo un mismo esquema, aunque no optimiza por separado la receta de cada red. Un estudio posterior podría comparar transferencia en ambas arquitecturas y establecer un presupuesto de búsqueda equivalente. Esas variantes requieren nuevos experimentos y un conjunto de prueba reservado para su evaluación.

Los giros amplios, volteos y relleno negro reproducen la intención de aumentos de la referencia con una implementación común. Pueden producir vistas de plausibilidad anatómica limitada. Un análisis adicional podría estudiar su efecto mediante validación y fijar una configuración alternativa antes de evaluar un nuevo conjunto de prueba. El protocolo actual permite atribuir sus resultados a decisiones específicas que están registradas en código.

La latencia se midió con entradas preparadas y lote uno. Un sistema que lea imágenes desde disco, reciba estudios completos o aplique otros preprocesamientos tendría un costo adicional. Los tiempos registrados permiten comparar las implementaciones probadas y deben acompañarse de la identificación de hardware y software. La cifra de entrenamiento incluye validación y guardado, por lo que su interpretación considera tanto las épocas completadas como la velocidad de cada paso.
